\section{Materials and Methods}
\label{sec:methods}

\subsection{Farthest-Point Landmark Sampling and Metric Basis Construction}
\label{subsec:methods_landmark_sampling}

Let $\mathcal{S} = \{s_1, s_2, \dots, s_M\}$ denote a multiple sequence alignment of $M$ contemporary nucleotide sequences of length $L$ over the alphabet $\Sigma = \{\text{A}, \text{C}, \text{G}, \text{T}\} \cup \{-\}$.
When $M \gg 1,000$, computing an all-to-all pairwise distance matrix requires $\binom{M}{2} \approx \frac{1}{2} M^2$ pairwise comparisons, scaling as $O(M^2 L)$ and creating a computational bottleneck for large-scale genomic surveillance.
Diversity anchors coordinates.
\RhizAeon eliminates this quadratic barrier through sequential Farthest-Point Sampling (FPS) \cite{gonzalez1985clustering}, selecting a sparse subset of $N \ll M$ reference landmark genomes in strictly $O(N M L)$ operations.

The selection algorithm begins by designating the initial reference landmark $s_{(1)}$ as the sequence exhibiting the highest non-gap coverage across the alignment:
\begin{equation}
s_{(1)} = \arg\max_{s_i \in \mathcal{S}} \sum_{u=1}^L \mathbb{I}(s_i(u) \neq \text{'-'}).
\label{eq:methods_fps_seed}
\end{equation}
The algorithm then computes the pairwise evolutionary distance from $s_{(1)}$ to all $M$ sequences in the cohort, initializing a running minimum distance vector $\mathbf{r} = (r_1, \dots, r_M)^T \in \mathbb{R}^M$, where $r_i = d(s_i, s_{(1)})$.
Pairwise evolutionary distances $d(s_i, s_j)$ are computed under the Tamura-Nei (TN93) continuous-time Markov substitution model with pairwise deletion of gap positions \cite{tamura1993estimation}:
\begin{equation}
d_{\text{TN93}} = -2 a_1 \ln\left(1 - \frac{P_1}{2 a_1} - \frac{Q}{2 a_3}\right) - 2 a_2 \ln\left(1 - \frac{P_2}{2 a_2} - \frac{Q}{2 a_3}\right) - 2 (a_3 - a_1 - a_2) \ln\left(1 - \frac{Q}{2 a_3}\right),
\label{eq:methods_tn93}
\end{equation}
where $P_1$ and $P_2$ represent transitional differences between purines and pyrimidines, $Q$ represents transversional differences, and $a_1, a_2, a_3$ denote equilibrium nucleotide frequency constants.
For alignments of highly divergent or non-coding sequences, normalized Hamming distance is supported as a scale-invariant alternative.

Each subsequent reference landmark $k \in \{2, \dots, N\}$ is selected sequentially by identifying the sequence that maximizes the minimum distance to the existing set of selected landmarks:
\begin{equation}
s_{(k)} = \arg\max_{i \in \{1,\dots,M\} \setminus \{s_{(1)},\dots,s_{(k-1)}\}} r_i.
\label{eq:methods_fps_iter}
\end{equation}
Upon selecting $s_{(k)}$, the algorithm evaluates distances from $s_{(k)}$ to all $M$ sequences and updates the running minimum distance array in $O(M L)$ time:
\begin{equation}
r_i \leftarrow \min\left(r_i, \, d(s_i, s_{(k)})\right), \quad \forall i \in \{1, \dots, M\}.
\label{eq:methods_fps_update}
\end{equation}
Selecting $N$ reference genomes across $M$ isolates requires exactly $N \times M$ distance evaluations, reducing the computational burden from $O(M^2 L)$ to $O(N M L)$.
For an alignment of 10,000 genomes with $N = 16$ landmarks, distance evaluations collapse from $49,995,000$ to $160,000$, representing a 312-fold empirical speedup.
Farthest points span divergence.

\paragraph{Full-Cohort Representation ($M \le 64$) and Two-Pass Clonal Backbone Re-Anchoring ($M > 64$).}
When scanning moderate surveillance cohorts ($M \le 64$), downsampling the alignment to a sparse subset of landmarks risks inadvertently selecting a recombinant isolate while omitting its parental lineage, or omitting a minor parental clade entirely.
To guarantee complete topological coverage, \RhizAeon automatically activates Full-Cohort Channel Tracking when $M \le 64$, setting the landmark set to the complete sequence roster ($N = M$).
Every sequence is endowed with its own continuous attention channel, eliminating downsampling artifacts and ensuring all prospective parental lineages are present in the reference basis.

For massive surveillance cohorts ($M > 64$), \RhizAeon utilizes Farthest-Point Sampling to select $N \le 32$ landmarks.
However, in epidemic recombinant swarms, greedy farthest-point exploration can pick an expanded recombinant lineage as a landmark due to its intermediate divergence.
If a recombinant isolate $R$ is admitted into the landmark basis, co-recombinant sister isolates ($R_1, \dots, R_k$) that share the identical mosaic cassettes will place dominant continuous attention ($\ge 80\%$) onto $R$.
When the parental resolver assigns Home = $R$, the differential excursion vanishes ($Y(u) \approx 0$), extinguishing the recombination signal in a catastrophic \textit{Landmark Self-Masking Trap}.
To eliminate this failure mode, \RhizAeon executes Two-Pass Clonal Backbone Re-Anchoring for $M > 64$:
Pass~1 screens the cohort using initial FPS landmarks to identify candidate recombinants ($\mathcal{C}_{\text{rec}}$).
If any landmark belongs to $\mathcal{C}_{\text{rec}}$, Pass~2 purges all candidate recombinants from the candidate landmark pool and re-anchors the landmark basis strictly onto verified non-recombinant clonal backbone lineages ($\mathcal{S} \setminus \mathcal{C}_{\text{rec}}$).
This guarantees that only stable, unadmixed parental lineages serve as coordinate landmarks.

From the selected reference genomes $\{s_{(1)}, \dots, s_{(N)}\}$, we construct the symmetric reference distance matrix $\mathbf{D} \in \mathbb{R}^{N \times N}$, with elements $D_{i,j} = d(s_{(i)}, s_{(j)})$ and diagonal entries $D_{i,i} = 0$.
We define the element-wise squared distance matrix $\mathbf{D}^{(2)} \in \mathbb{R}^{N \times N}$ such that $[\mathbf{D}^{(2)}]_{i,j} = D_{i,j}^2$.
To map these distances into an inner-product space, we apply the geometric centering projection operator $\mathbf{H} = \mathbf{I}_N - \frac{1}{N} \mathbf{1}_N \mathbf{1}_N^T$, where $\mathbf{I}_N$ is the $N \times N$ identity matrix and $\mathbf{1}_N = (1, \dots, 1)^T \in \mathbb{R}^N$ is the all-ones vector.
The doubly centered inner-product Gram matrix $\mathbf{B} \in \mathbb{R}^{N \times N}$ is formulated via classical multidimensional scaling (MDS) \cite{gower1966some}:
\begin{equation}
\mathbf{B} = -\frac{1}{2} \mathbf{H} \mathbf{D}^{(2)} \mathbf{H}.
\label{eq:methods_gram_matrix}
\end{equation}
Spectral decomposition of the symmetric matrix $\mathbf{B}$ yields:
\begin{equation}
\mathbf{B} = \mathbf{V} \mathbf{\Lambda} \mathbf{V}^T = \sum_{k=1}^N \lambda_k \mathbf{v}_k \mathbf{v}_k^T,
\label{eq:methods_spectral_decomp}
\end{equation}
where $\mathbf{\Lambda} = \operatorname{diag}(\lambda_1, \lambda_2, \dots, \lambda_N)$ contains real eigenvalues arranged in descending order ($\lambda_1 \ge \lambda_2 \ge \dots \ge \lambda_N$), and $\mathbf{V} = [\mathbf{v}_1, \dots, \mathbf{v}_N] \in \mathbb{R}^{N \times N}$ is the orthonormal eigenvector matrix satisfying $\mathbf{V}^T \mathbf{V} = \mathbf{I}_N$.
Spectra reveal splits.

We retain the leading $d$ dimensions corresponding to the largest positive eigenvalues ($d \le \min(4, N-1)$):
\begin{equation}
\mathbf{\Lambda}_d = \operatorname{diag}(\lambda_1, \dots, \lambda_d) \in \mathbb{R}^{d \times d}, \quad \mathbf{V}_d = [\mathbf{v}_1, \dots, \mathbf{v}_d] \in \mathbb{R}^{N \times d}.
\label{eq:methods_retained_eigen}
\end{equation}
The static $d$-dimensional metric coordinates $\mathbf{x}_i \in \mathbb{R}^d$ for reference landmark $i \in \{1, \dots, N\}$ are given by:
\begin{equation}
\mathbf{x}_i = \mathbf{\Lambda}_d^{1/2} (\mathbf{V}_d)_{i, :}^T = \left(\sqrt{\lambda_1} v_{1, i}, \, \sqrt{\lambda_2} v_{2, i}, \, \dots, \, \sqrt{\lambda_d} v_{d, i}\right)^T.
\label{eq:methods_static_coords}
\end{equation}
By construction, these coordinates satisfy $\langle \mathbf{x}_i, \mathbf{x}_j \rangle = [\mathbf{B}]_{i,j}$ within the subspace $\mathbb{R}^d$.
By Buneman's metric tree theorem \cite{buneman1974note}, tree metrics satisfy four-point additivity conditions $d(a, b) + d(c, d) \le \max(d(a, c) + d(b, d), \, d(a, d) + d(b, c))$, ensuring that tree-like evolutionary divergences embed into Euclidean spaces where macroscopic clade splits project onto the dominant eigenvectors.

\subsection{Analytical Gower Out-of-Sample Trajectory Projection}
\label{subsec:methods_gower_projection}

Tracking sequence affinities along the chromosome requires establishing an instantaneous coordinate $\mathbf{z}_i(u) \in \mathbb{R}^d$ for each query sequence $i \in \{1, \dots, M\}$ at every genomic site $u \in [1, L]$.
Recomputing multidimensional scaling independently within sliding windows would produce arbitrary orthogonal rotations and reflections, distorting geometric distances and necessitating expensive Procrustes alignments.
Instead, \RhizAeon projects local distance vectors into the fixed canonical coordinate basis $(\mathbf{V}_d, \mathbf{\Lambda}_d)$ using Gower's analytical out-of-sample interpolation theorem \cite{gower1968adding}.

At coordinate $u$, we evaluate local evolutionary distances across an aligned window $[u - W/2, u + W/2]$ of width $W$ (default $W = 300\text{ nt}$) sampled at stride $S$ (default $S = 10\text{ nt}$).
For query sequence $i$, this yields an instantaneous local distance vector $\mathbf{d}_i(u) \in \mathbb{R}^N$ to the $N$ reference landmarks:
\begin{equation}
\mathbf{d}_i(u) = \left(d_i(1; u), \, d_i(2; u), \, \dots, \, d_i(N; u)\right)^T.
\label{eq:methods_local_dist_vec}
\end{equation}
We define the vector of squared local distances as $\mathbf{d}_i^2(u) = (d_i(1; u)^2, \dots, d_i(N; u)^2)^T \in \mathbb{R}^N$, and the vector of mean squared reference distances as:
\begin{equation}
\bar{\mathbf{d}}^2 = \frac{1}{N} \mathbf{D}^{(2)} \mathbf{1}_N = \left( \frac{1}{N} \sum_{j=1}^N D_{1, j}^2, \, \dots, \, \frac{1}{N} \sum_{j=1}^N D_{N, j}^2 \right)^T \in \mathbb{R}^N.
\label{eq:methods_mean_sq_dist}
\end{equation}

Gower's out-of-sample mapping identifies the coordinate $\mathbf{z}_i(u) \in \mathbb{R}^d$ that minimizes squared Euclidean distance distortions relative to the static reference coordinates $\mathbf{x}_j$:
\begin{equation}
\min_{\mathbf{z} \in \mathbb{R}^d} \sum_{j=1}^N \left( \|\mathbf{z} - \mathbf{x}_j\|^2 - d_i(j; u)^2 \right)^2.
\label{eq:methods_gower_objective}
\end{equation}
Expanding the Euclidean norm $\|\mathbf{z} - \mathbf{x}_j\|^2 = \|\mathbf{z}\|^2 - 2 \mathbf{z}^T \mathbf{x}_j + \|\mathbf{x}_j\|^2$, utilizing the centering condition $\sum_{j=1}^N \mathbf{x}_j = \mathbf{0}$, and equating the gradient to zero yields the exact analytical solution:
\begin{equation}
\mathbf{z}_i(u) = \frac{1}{2} \mathbf{\Lambda}_d^{-1/2} \mathbf{V}_d^T \left( \bar{\mathbf{d}}^2 - \mathbf{d}_i^2(u) \right) \in \mathbb{R}^d,
\label{eq:methods_gower_formula}
\end{equation}
where $\mathbf{\Lambda}_d^{-1/2} = \operatorname{diag}(1/\sqrt{\lambda_1}, \dots, 1/\sqrt{\lambda_d})$.
Evaluating \eqref{eq:methods_gower_formula} across all coordinates $u \in \{1, 1+S, \dots, L\}$ constructs the continuous trajectory tensor $\mathbf{Z} \in \mathbb{R}^{M \times U \times d}$, where $U = \lceil L/S \rceil$.

Because the basis $(\mathbf{V}_d, \mathbf{\Lambda}_d)$ and mean vector $\bar{\mathbf{d}}^2$ are computed once from whole-genome reference divergence, every sliding window maps into the identical, rigid metric reference frame.
Zero Procrustes rotations are required.
Projections preserve geometry.
Furthermore, substituting the whole-genome distance vector of reference landmark $k$ into \eqref{eq:methods_gower_formula} analytically recovers its exact MDS coordinate:
\begin{equation}
\frac{1}{2} \mathbf{\Lambda}_d^{-1/2} \mathbf{V}_d^T \left( \bar{\mathbf{d}}^2 - \left( \mathbf{D}^{(2)}_{k, :} \right)^T \right) = \mathbf{\Lambda}_d^{1/2} (\mathbf{V}_d)_{k, :}^T = \mathbf{x}_k,
\label{eq:methods_gower_identity}
\end{equation}
proving algebraic self-consistency and zero metric distortion on reference taxa.

\subsection{Generalized Tree-RoPE and Continuous-Time Markov Substitution Prior}
\label{subsec:methods_treerope_prior}

Standard transformer neural networks incorporate sequential position along one-dimensional text using Rotary Position Embeddings (RoPE) \cite{su2024roformer}.
However, biological sequences diverge across multidimensional phylogenetic trees rather than linear strings.
To integrate evolutionary distance directly into cross-sequence attention, \RhizAeon generalizes rotary embeddings from one-dimensional sequence indices to continuous metric coordinates.

Let $\mathbf{x}_i \in \mathbb{R}^d$ denote the continuous metric coordinates of sequence $i$.
We construct a fixed frequency projection basis $\mathbf{\Theta} \in \mathbb{R}^{(d_{\text{head}}/2) \times d}$, with components parameterized by base frequency $\theta_0 = 10,000$:
\begin{equation}
\Theta_{m, c} = \theta_0^{-2(m-1)/d_{\text{head}}} \cdot \omega_c, \quad m \in \{1, \dots, d_{\text{head}}/2\}, \; c \in \{1, \dots, d\},
\label{eq:methods_freq_basis}
\end{equation}
where $\omega_c = \sqrt{\lambda_c / \sum_{k=1}^d \lambda_k}$ weights each dimension by its relative spectral variance.
For sequence $i$, the rotary phase angle vector $\bm{\theta}_i = (\theta_{i, 1}, \dots, \theta_{i, d_{\text{head}}/2})^T$ is computed as:
\begin{equation}
\theta_{i, m} = \sum_{c=1}^d x_{i, c} \Theta_{m, c}.
\label{eq:methods_tree_rope_angle}
\end{equation}
Angles encode divergence.

At each aligned site $u$, sequence representations generate query and key vectors $\mathbf{q}_i(u), \mathbf{k}_i(u) \in \mathbb{R}^{d_{\text{head}}}$.
We define the block-diagonal Givens rotation matrix $\mathbf{R}(\bm{\theta}_i) \in \mathbb{R}^{d_{\text{head}} \times d_{\text{head}}}$ as:
\begin{equation}
\mathbf{R}(\bm{\theta}_i) = \bigoplus_{m=1}^{d_{\text{head}}/2} \begin{pmatrix} \cos \theta_{i, m} & -\sin \theta_{i, m} \\ \sin \theta_{i, m} & \cos \theta_{i, m} \end{pmatrix}.
\label{eq:methods_givens_rot}
\end{equation}
Applying $\mathbf{R}(\bm{\theta}_i)$ rotates query and key representations according to each taxon's continuous metric position.
The inner product between rotated vectors satisfies:
\begin{equation}
\langle \mathbf{R}(\bm{\theta}_i) \mathbf{q}_i(u), \, \mathbf{R}(\bm{\theta}_j) \mathbf{k}_j(u) \rangle = \mathbf{q}_i(u)^T \mathbf{R}(\bm{\theta}_j - \bm{\theta}_i) \mathbf{k}_j(u).
\label{eq:methods_rot_inner}
\end{equation}
Because rotation angles depend linearly on metric coordinates, the phase shift $\|\bm{\theta}_j - \bm{\theta}_i\|$ increases with evolutionary divergence in Buneman metric space, naturally attenuating pairwise attention as phylogenetic distance increases.

To calibrate geometric distances against baseline neutral sequence homology, we integrate a continuous-time Markov substitution prior directly into the attention logits.
For any pair of sequences $(i, j)$ separated by evolutionary distance $D_{i,j}$, the probability of observing identical homologous nucleotides under a continuous-time Markov substitution model with substitution scaling rate $\lambda_h$ and equilibrium frequency $\epsilon_0 = 0.25$ is:
\begin{equation}
P_{i, j} = \epsilon_0 + (1 - \epsilon_0) \exp\left(-\lambda_h D_{i, j}\right).
\label{eq:methods_markov_prior}
\end{equation}
We define the logarithmic phylogenetic prior bias as:
\begin{equation}
\text{PhyloBias}_{i, j} = \ln \left[ \max\left(10^{-5}, \, P_{i, j}\right) \right].
\label{eq:methods_phylobias}
\end{equation}
The composite unnormalized cross-attention score $S_u(i, j)$ between query sequence $i$ and candidate reference $j$ at coordinate $u$ becomes:
\begin{equation}
S_u(i, j) = \frac{(\mathbf{R}(\bm{\theta}_i) \mathbf{q}_i(u))^T (\mathbf{R}(\bm{\theta}_j) \mathbf{k}_j(u))}{\sqrt{d_{\text{head}}}} + \text{PhyloBias}_{i, j}.
\label{eq:methods_composite_score}
\end{equation}

\subsection{Unpolarized Root Sink Architecture and Autapomorphy Absorption}
\label{subsec:methods_root_sink}

Standard cross-attention mechanisms normalize probability exclusively across contemporary leaves, creating structural vulnerability to lineage-specific rate variation.
Private mutations lack donors.
When an isolated lineage undergoes an episodic substitution burst---such as APOBEC3G cytidine deamination or accelerated heterotachy---private mutations lack homologous counterparts in contemporary clades.
Without an unpolarized reference, softmax forces these private mutations to affiliate with whichever contemporary leaf presents the smallest stochastic mismatch, generating severe false positive recombination calls.

To insulate trajectory inference from unreticulated rate bursts, \RhizAeon reserves index zero in the sequence dimension for an unpolarized ancestral root token $\mathbf{x}_{\text{root}}$.
The root token is anchored at the origin of the Buneman metric space:
\begin{equation}
\mathbf{x}_{\text{root}} = \mathbf{0} = (0, 0, \dots, 0)^T \in \mathbb{R}^d.
\label{eq:methods_root_anchor}
\end{equation}
For an alignment of $N$ contemporary reference genomes, the cross-attention operator operates over $N+1$ nodes, spanning the ancestral root ($j=0$) and all contemporary leaves ($j=1, \dots, N$).
The softmax normalization strictly conserves probability across all $N+1$ nodes:
\begin{equation}
A_u(i \to j) = \frac{\exp(S_u(i, j))}{\sum_{k=0}^N \exp(S_u(i, k))}, \quad \text{such that} \quad A_u(i \to 0) + \sum_{j=1}^N A_u(i \to j) = 1.0.
\label{eq:methods_root_softmax}
\end{equation}

\begin{proof}[Proof of Autapomorphy Absorption]
Consider a genomic site $u$ where sequence $i$ acquires a private autapomorphic mutation that does not match any contemporary reference genome $j \in \{1, \dots, N\}$.
For all contemporary leaves $j \ge 1$, the mismatch depresses both query-key matching and the phylogenetic prior, yielding low scores $S_u(i, j) \le S_{\text{mismatch}}$.
In contrast, the unpolarized root token at $\mathbf{x}_{\text{root}} = \mathbf{0}$ represents the baseline unmutated ancestral state with neutral substitution distance, maintaining $S_u(i, 0) = S_{\text{root}} > S_{\text{mismatch}}$.
In the limit where private divergence increases:
\begin{equation}
\lim_{S_u(i, j) \to -\infty, \, \forall j \ge 1} A_u(i \to 0) = \frac{\exp(S_u(i, 0))}{\exp(S_u(i, 0)) + \sum_{j=1}^N \exp(S_u(i, j))} = 1.0,
\end{equation}
and $A_u(i \to j) \to 0.0$ for all $j \ge 1$.
Consequently, autapomorphic substitution bursts route their attention flux into node zero rather than contemporary leaves, preventing spurious donor excursions and guaranteeing a 0.00\% false positive rate under biological heterotachy.
\end{proof}
Sinks absorb outliers.

\subsection{Continuous Trajectory Integration and Brownian Bridge Inference}
\label{subsec:methods_brownian_bridge}

At individual polymorphic sites, instantaneous attention weights exhibit stochastic high-frequency jitter.
To extract continuous evolutionary trajectories, \RhizAeon integrates attention into cumulative trajectory curves.
Let $\mathcal{B}_m \subset \{1, \dots, N\}$ denote the set of reference genomes belonging to phylogenetic clade $m$.
The instantaneous bundle attention is $\alpha_{i, m}(u) = \sum_{j \in \mathcal{B}_m} A_u(i \to j)$.
The continuous cumulative attention curve $C_{i, m}(u)$ represents the integrated affinity to clade $m$ up to coordinate $u$:
\begin{equation}
C_{i, m}(u) = \sum_{v=1}^u \alpha_{i, m}(v) = \sum_{v=1}^u \sum_{j \in \mathcal{B}_m} A_v(i \to j).
\label{eq:methods_cum_curve}
\end{equation}

The home clade $m_{\text{home}}(i)$ is identified as the lineage accumulating the greatest integrated affinity across the full chromosome:
\begin{equation}
m_{\text{home}}(i) = \arg\max_m C_{i, m}(L).
\label{eq:methods_home_clade}
\end{equation}
For any candidate donor clade $m \neq m_{\text{home}}$, we formulate the differential attention trajectory as:
\begin{equation}
Y_{i, m}(u) = C_{i, m}(u) - C_{i, m_{\text{home}}}(u).
\label{eq:methods_diff_trajectory}
\end{equation}
Under clonal evolution within the parental background, attention to the home clade exceeds donor attention at nearly all sites ($A_v(i \to \text{home}) > A_v(i \to m)$), causing $Y_{i, m}(u)$ to drift monotonically downward with negative slope.
Recombination reverses the derivative: within a converted donor tract $[u_1, u_2]$, donor attention dominates, producing a positive slope ($dY/du > 0$) and a prominent upward excursion.

To test the clonal null hypothesis $H_0$, we construct the centered Brownian bridge process $B_{i, m}(u)$ by subtracting the chromosome-wide linear drift:
\begin{equation}
B_{i, m}(u) = Y_{i, m}(u) - \frac{u}{L} Y_{i, m}(L), \quad u \in [1, L].
\label{eq:methods_brownian_bridge}
\end{equation}
Let $\Delta Y(u) = Y(u) - Y(u-1)$ denote the local trajectory increments, with sample mean $\bar{\Delta Y} = \frac{1}{L} Y(L)$ and empirical increment variance:
\begin{equation}
\sigma^2 = \frac{1}{L - 1} \sum_{u=1}^L \left( \Delta Y(u) - \bar{\Delta Y} \right)^2.
\label{eq:methods_sigma_sq}
\end{equation}

Under the clonal null hypothesis $H_0$, the increments are stationary and unreticulated.
By Donsker's functional central limit theorem \cite{donsker1951invariance}, the standardized process:
\begin{equation}
W_L^0(t) = \frac{B_{i, m}(\lfloor t L \rfloor)}{\sigma \sqrt{L}}, \quad t \in [0, 1],
\label{eq:methods_donsker_bridge}
\end{equation}
converges weakly as $L \to \infty$ to a standard Brownian bridge $W^0(t)$ with continuous sample paths, mean $\mathbb{E}[W^0(t)] = 0$, and covariance $\operatorname{Cov}(W^0(s), W^0(t)) = \min(s, t) - s t$.
The supremum of the absolute bridge trajectory follows the two-sided Kolmogorov distribution:
\begin{equation}
P\left( \sup_{t \in [0, 1]} |W^0(t)| \le z \right) = 1 - 2 \sum_{k=1}^\infty (-1)^{k-1} \exp\left(-2 k^2 z^2\right).
\label{eq:methods_kolmogorov_cdf}
\end{equation}

To control the family-wise error rate (FWER) across all $M_{\text{topo}} = \binom{N}{2} \max(1, N - 2)$ testable donor-recipient reticulation topologies at nominal significance $\alpha = 0.05$, we apply the Dunn-\v{S}id\'{a}k extreme-value correction \cite{sidak1967rectangular}:
\begin{equation}
K_{\text{crit}} = \sqrt{-\frac{1}{2} \ln \left( 1 - (1 - \alpha)^{1 / M_{\text{topo}}} \right)} \approx \sqrt{-\frac{1}{2} \ln \left( \frac{\alpha}{2 M_{\text{topo}}} \right)}.
\label{eq:methods_sidak_bound}
\end{equation}
If the normalized trajectory supremum satisfies:
\begin{equation}
\frac{\sup_{u \in [1, L]} |B_{i, m}(u)|}{\sigma \sqrt{L}} \ge K_{\text{crit}},
\label{eq:methods_rejection_rule}
\end{equation}
\RhizAeon rejects the clonal null hypothesis and confirms a statistically significant reticulation event.

\subsection{Dyadic Scale-Space Matched Filtering and Running Drawdown}
\label{subsec:methods_scalespace_drawdown}

For short genetic conversions and micro-recombinations ($< 200\text{ nt}$), global Brownian bridge tests can suffer signal dilution across multi-kilobase chromosomes.
To achieve multi-scale sensitivity, \RhizAeon convolves the differential trajectory flux $\Delta Y(u)$ with a bank of zero-mean bipolar dyadic matched filters $h_\tau(s)$ evaluated across scales $\tau \in \{15, 25, 50, 100, 200, 500\}$ nucleotides:
\begin{equation}
h_\tau(s) = \begin{cases} +1, & 0 \le s < \tau/2 \\ -1, & \tau/2 \le s < \tau \\ 0, & \text{otherwise} \end{cases}.
\label{eq:methods_wavelet}
\end{equation}
The discrete convolution $(dY * h_\tau)(u) = \sum_{s=0}^\tau \Delta Y(u - s) h_\tau(s)$ isolates rapid step transitions across dyadic scale space.

Concurrently, \RhizAeon tracks the running high-water mark $M(u)$ and running drawdown $\text{DD}(u)$:
\begin{equation}
M(u) = \max_{1 \le v \le u} Y(v), \quad \text{DD}(u) = M(u) - Y(u).
\label{eq:methods_running_dd}
\end{equation}
The drawdown metric localizes points where trajectory momentum collapses, marking the exit boundary $u_2$ of a donor tract.
To distinguish single crossovers from intercalary cassette swaps, we evaluate the normalized energy density ratio of the trajectory excursions outside versus inside the candidate tract $[u_1, u_2]$:
\begin{equation}
\mathcal{E}_{\text{ratio}} = \frac{\frac{1}{L - (u_2 - u_1)} \int_{u \notin [u_1, u_2]} (Y(u) - Y_{\text{base}})^2 \, du}{\frac{1}{u_2 - u_1} \int_{u_1}^{u_2} (Y(u) - Y_{\text{base}})^2 \, du},
\label{eq:methods_energy_ratio}
\end{equation}
where $Y_{\text{base}}$ denotes the linear baseline drift.
A single crossover spans half the chromosome, producing high residual energy across the alignment ($\mathcal{E}_{\text{ratio}} \ge 0.50$).
In contrast, an intercalary cassette returns cleanly to baseline on both flanks, characterized by $\mathcal{E}_{\text{ratio}} < 0.15$.
Complementing the trajectory energy ratio, for decomposed spatial modes $\mathbf{u}_m$, \RhizAeon computes the continuous spatial mode polarity ratio:
\begin{equation}
p_{\text{mono}} = \frac{\left| \sum_{u=1}^L u_m(u) \right|}{\sum_{u=1}^L |u_m(u)|}.
\label{eq:methods_polarity_ratio}
\end{equation}
A monopolar arch maintaining uniform sign across the alignment yields $p_{\text{mono}} \approx 1.0$ (characterizing single crossovers and tree backbones), whereas a localized bipolar pulse returning to baseline on both flanks yields $p_{\text{mono}} \le 0.50$ (characterizing intercalary cassette swaps; Figure~\ref{fig:theory_force_field}D).

\subsection{Physical Force Field Deconvolution and Random Matrix Halting}
\label{subsec:methods_force_field}

To deconvolve complex multi-way mosaicism, \RhizAeon models trajectory excursions as physical displacements within an evolutionary force field.
We parameterize the evolutionary pull of the home lineage as a Hookean restoring potential $U(\mathbf{z}) = \frac{1}{2} \lambda \|\mathbf{z} - \mathbf{z}_{\text{home}}\|^2$, generating a continuous restoring force:
\begin{equation}
\mathbf{f}(u) = -\nabla U = -\lambda \left(\mathbf{z}(u) - \mathbf{z}_{\text{home}}\right) \in \mathbb{R}^d.
\label{eq:methods_force_vector}
\end{equation}

Across a candidate breakpoint boundary, let $\mathbf{f}_{\text{left}}$ and $\mathbf{f}_{\text{right}}$ denote the mean force vectors on the flanking intervals.
The directional deflection angle $\theta$ satisfies $\cos \theta = \frac{\langle \mathbf{f}_{\text{left}}, \mathbf{f}_{\text{right}} \rangle}{\|\mathbf{f}_{\text{left}}\| \|\mathbf{f}_{\text{right}}\|}$, and the step velocity is $v_{\text{step}} = \|\mathbf{f}_{\text{right}} - \mathbf{f}_{\text{left}}\|$.
The dynamic velocity reversal score is defined as:
\begin{equation}
R = v_{\text{step}} (1 - \cos \theta).
\label{eq:methods_velocity_reversal}
\end{equation}
To achieve scale-invariance across cohorts of varying divergence, we standardize the velocity reversal along the unit parental dipole axis $\mathbf{d}_{\text{hat}} = \frac{\mathbf{z}_{\text{donor}} - \mathbf{z}_{\text{home}}}{\|\mathbf{z}_{\text{donor}} - \mathbf{z}_{\text{home}}\|}$:
\begin{equation}
Z_{\text{flip}} = \frac{\left| (\mathbf{f}_{\text{right}} - \mathbf{f}_{\text{left}}) \cdot \mathbf{d}_{\text{hat}} \right|}{\sqrt{\frac{\sigma_{\text{left}}^2}{L_{\text{left}}} + \frac{\sigma_{\text{right}}^2}{L_{\text{right}}}}},
\label{eq:methods_z_flip}
\end{equation}
where $\sigma_{\text{left}}^2, \sigma_{\text{right}}^2$ represent residual force variances along $\mathbf{d}_{\text{hat}}$ across flanking segments of lengths $L_{\text{left}}, L_{\text{right}}$.

For alignments exhibiting multiple recombinant tracts, we stack force vectors into the spatio-temporal force matrix $\mathbf{F} \in \mathbb{R}^{N \times L}$ and perform singular value decomposition:
\begin{equation}
\mathbf{F} = \sum_{m=1}^r s_m \mathbf{u}_m \mathbf{v}_m^T,
\label{eq:methods_svd_peeling}
\end{equation}
where $s_1 \ge s_2 \ge \dots \ge s_r$ denote singular values.
Each detected reticulation mode $m$ is deconvolved via iterative mode peeling, subtracting $s_m \mathbf{u}_m \mathbf{v}_m^T$ from $\mathbf{F}$.
To prevent over-segmentation and halt peeling when residual variation reflects neutral coalescent drift, we apply random matrix theory.
Under the unreticulated null hypothesis, the singular values of an $N \times L$ random noise matrix follow the Marchenko-Pastur distribution \cite{marchenko1967distribution}.
The theoretical upper boundary of the bulk spectral noise edge is:
\begin{equation}
s_{\text{bulk}} = \frac{L}{\pi} \sigma_{\text{noise}} \left( 1 + \sqrt{\frac{P}{L}} \right),
\label{eq:methods_marchenko_pastur}
\end{equation}
where $P = \min(N, d)$ denotes the metric projection dimension and $\sigma_{\text{noise}}$ represents the baseline standard deviation of empirical force increments.
Mode peeling halts when the leading residual singular value satisfies $s_m \le s_{\text{bulk}}$, terminating iteration without manual parameter tuning.

\subsection{Profile Maximum Likelihood Breakpoint Polishing}
\label{subsec:methods_polisher}

Following changepoint localization from continuous trajectory dynamics, \RhizAeon refines candidate breakpoints down to single-nucleotide resolution using profile maximum likelihood polishing.
Within a local search radius $[u_{\text{init}} - \delta, u_{\text{init}} + \delta]$ (default $\delta = 50\text{ nt}$), the engine scans each candidate coordinate $c$ and computes the partitioned log-likelihood:
\begin{equation}
\ln \mathcal{L}(c) = \sum_{u=1}^c \ln P\left(s_{\text{rec}}(u) \mid s_{\text{left}}(u)\right) + \sum_{u=c+1}^L \ln P\left(s_{\text{rec}}(u) \mid s_{\text{right}}(u)\right),
\label{eq:methods_profile_likelihood}
\end{equation}
where $P(s_{\text{rec}}(u) \mid s(u))$ denotes the conditional probability of the recombinant state given the parental sequence under the fitted nucleotide substitution model.

When invariant or uninformative sites intervene between flanking parental polymorphisms, the likelihood function forms a flat plateau.
Rather than reporting an arbitrary midpoint, \RhizAeon reports the exact informational plateau $[c_{\text{left}}, c_{\text{right}}]$:
\begin{equation}
[c_{\text{left}}, c_{\text{right}}] = \left\{ c \in [1, L] : \ln \mathcal{L}(c) \ge \max_{c'} \ln \mathcal{L}(c') - \Delta_{\text{thresh}} \right\},
\label{eq:methods_info_plateau}
\end{equation}
where $\Delta_{\text{thresh}} = 1.92$ corresponds to the 95\% confidence interval under a $\chi_1^2$ likelihood ratio test.
The boundary coordinates $c_{\text{left}}$ and $c_{\text{right}}$ demarcate the exact flanking informative polymorphisms that physically bound the genetic exchange.
Plateaus bound breakpoints.

\paragraph{Continuous Gap Bridging Across Stochastic Non-Informative Dips.}
In low-divergence regimes ($d \le 0.02$), mutation density is inherently sparse ($\sim 1$ substitution every 50 nucleotides).
Under standard Poisson mutation processes, stochastic intervals of 50--100 bp within an authentic 1,000 nt conversion cassette frequently carry zero donor-specific mutations ($P(0 \text{ mutations}) = e^{-\lambda} \approx 15\%\text{--}35\%$).
During recursive binary segmentation, instantaneous differential flux drops to zero across these unmutated windows, which would otherwise fracture an intact biological cassette into multiple fragmented micro-tracts.
To preserve macro-conversion integrity, \RhizAeon implements Continuous Gap Bridging:
adjacent candidate donor segments separated by small non-informative gaps ($\le 150$ bp) are iteratively coalesced into a single unified tract whenever the net flux across the merged span remains positive ($\frac{1}{e_2 - s_1 + 1} \sum_{u=s_1}^{e_2} X(u) > 0$).
Bridging restores the contiguous physical envelope of the biological introgression, pooling all supporting mutations together and ensuring the subsequent hypergeometric contingency test operates on the true cumulative evidence.

\subsection{Reticulation Graph Adjudication, Transitive Chains, and Ancestral Event Grouping}
\label{subsec:methods_adjudication_grouping}

In pairwise and triplet sequence scans, an intrinsic topological ambiguity arises: if a recombinant isolate $C$ inherits a genomic tract from donor $D$ into a parental home background $H$, reciprocal tests frequently evaluate $D$ as recombinant between $C$ and $H$.
This reciprocal false alarm creates a mutual two-cycle in candidate reticulation graphs ($C \leftrightarrow D$).
To orient parental lineages versus genuine recombinant isolates without combinatorial tree searches, \RhizAeon inspects candidate events across homologous intervals.
Homologous candidates share overlapping tracts ($\ge 25\%$ of the shorter tract) or form complementary crossover partitions ($u_2^{(1)} \approx u_1^{(2)}$).

The algorithm evaluates external donor authenticity.
If candidate $C_1$ identifies an external donor while candidate $C_2$ only cites $C_1$, the algorithm prunes $C_2$ as the non-recombinant parental donor.
When candidates cite each other reciprocally, \RhizAeon computes physical transition contrast and home backbone disruption across the breakpoint boundaries.
The contrast metrics measure differential divergence inside versus outside the candidate tract:
\begin{align}
\Delta d_{\text{home}} &= \max\left(0, \, d(C, H \mid \text{tract}) - d(C, H \mid \text{flank})\right), \label{eq:methods_delta_home} \\
\Delta d_{\text{donor}} &= \max\left(0, \, d(C, D \mid \text{flank}) - d(C, D \mid \text{tract})\right). \label{eq:methods_delta_donor}
\end{align}
A true recombinant child experiences substantial home backbone disruption ($\Delta d_{\text{home}} \gg 0$) by departing from its parental home basin during the conversion tract.
In contrast, a stationary clonal parent remains within its home background throughout the alignment ($\Delta d_{\text{home}} \approx 0$).
Comparing total transition contrast $\Delta_{\text{trans}} = \Delta d_{\text{home}} + \Delta d_{\text{donor}}$ and mosaic reconstruction residual $d(C, D \mid \text{tract}) + d(C, H \mid \text{flank})$ unambiguously orients the recombinant isolate and purges reciprocal parent false alarms.

\paragraph{Intra-Swarm Sister Isolation, Donor Clade Deduplication, and Sister Preservation.}
In expanding epidemic swarms, multiple recombinant isolates descend from a common recombinant ancestor, differing only by minor intra-swarm drift ($d < 0.010$).
Without explicit sister safeguards, three distinct failure modes corrupt swarm deconvolution:
First, during parental channel resolution, co-recombinant sisters place immense continuous attention onto each other; \RhizAeon strictly excludes contemporary channels within intra-swarm drift ($d(s_i, s_j) < 0.010$) from serving as prospective Home or Donor channels for each other, forcing attention to resolve against genuine parental clades.
Second, in multi-way mosaics acquiring distinct cassettes from separate donor clades (e.g., $P_2$ and $P_3$), candidate donors from the primary clade ($P_2$) can monopolize the candidate pool; \RhizAeon deduplicates candidate donors by phylogenetic distance ($d(d_1, d_2) < 0.010$), guaranteeing that secondary donor clades are not crowded out by redundant primary sisters.
Third, in reticulation graph adjudication, reciprocal two-cycle pruning must not confuse parallel sister isolates with reciprocal parental illusions; \RhizAeon explicitly protects co-recombinant sister events sharing identical Home and Donor assignments ($\text{Home}_1 = \text{Home}_2$ and $\text{Donor}_1 = \text{Donor}_2$) from mutual reflection pruning.
All verified sister isolates are preserved and subsequently aggregated into the ancestral event roster $\mathcal{I} = \{s_1, \dots, s_k\}$.

In surveillance cohorts containing multiple related recombinants, an isolate $A$ may acquire a tract from an ancestral recombinant $B$ before $B$ further diverged or was sampled.
Under these conditions, isolate $A$ cites $B$ as its donor or home.
If $B$ was itself introgressed from donor $C$, this dependency creates a transitive donor chain ($A \to B \to C$).
\RhizAeon iteratively traverses donor and home dependency chains across overlapping tracts sharing $\ge 50\%$ tract overlap.
When isolate $A$ cites $B$ as donor, and $B$ (or its co-recombinant cluster) cites external donor $C$ with a compatible home lineage, the algorithm reassigns $A$ donor to $C$, updating tract and flanking divergences $d(A, C \mid \text{tract})$ and $d(A, C \mid \text{flank})$.
Symmetrically, if $A$ cites $B$ as home and $B$ cites $C$ as home with a compatible donor, $A$ inherits home $C$.
The chain traversal forbids cyclic attributions ($A \to B \to A$), resolving transitive co-recombinant cascades into primary evolutionary transfers.

When an ancestral recombination event precedes clonal expansion, multiple descendant isolates inherit the identical mosaic tract.
Reporting each isolate as an independent recombination event inflates event counts and misrepresents epidemic transmission.
\RhizAeon clusters verified events into shared ancestral recombination events using an interval overlap criterion.
Events sharing identical resolved Home and Donor lineages whose breakpoint intervals overlap by $\ge 50\%$ of the shorter tract:
\begin{equation}
\frac{\min(u_2^{(1)}, u_2^{(2)}) - \max(u_1^{(1)}, u_1^{(2)}) + 1}{\min(u_2^{(1)} - u_1^{(1)} + 1, \, u_2^{(2)} - u_1^{(2)} + 1)} \ge 0.50,
\label{eq:methods_ancestral_overlap}
\end{equation}
are collapsed into a single ancestral event.
The grouped event aggregates the complete roster of co-recombinant descendant isolates ($\mathcal{I} = \{s_1, \dots, s_k\}$).
The consensus breakpoint interval spans the outer envelope $[ \min_k u_{1, k}, \, \max_k u_{2, k} ]$, while preserving the maximal statistical support across member isolates ($\min p_{\text{Fisher}}$, $\max Z_{\text{phys}}$, $\max S_{\text{informative}}$).
Each event receives a persistent sequential identifier ($1, 2, \dots$), reporting both the count of unique evolutionary events and total recombinant isolates across the alignment inventory.

\subsection{Outgroup-Invariant Statistical Due Diligence}
\label{subsec:methods_due_diligence}

In cohorts with $N \ge 4$ taxa, \RhizAeon searches for an unrooted quartet outgroup $O$ that maximizes divergence to candidate and home lineages ($d(O, H) + d(O, D)$), preferentially selecting non-candidate, non-rate-burst background leaves to prevent outgroup poisoning.
If all external lineages have been admitted as candidates due to high sensitivity in low-divergence regimes ($d \le 0.02$), the selector falls back to the most divergent available taxon outside $\{C, H, D\}$, ensuring quartet verification never starves.
When an authentic external outgroup exists, the engine tests the unrooted Buneman four-point condition across tract and flank partitions:
\begin{align}
S_1 &= d(C, H) + d(D, O), \label{eq:methods_buneman_s1} \\
S_2 &= d(C, D) + d(H, O), \label{eq:methods_buneman_s2} \\
S_3 &= d(C, O) + d(H, D). \label{eq:methods_buneman_s3}
\end{align}
Under genuine recombination, the minimal parsimony split flips from $S_1$ (grouping $C$ with $H$ on flanks) to $S_2$ (grouping $C$ with $D$ in tract), with the topological difference exceeding the neutral binomial standard error ($z \ge z_{\text{crit}}$).

In two-clade epidemic expansions where all sampled sequences belong to parental clades or recombinant progeny, no authentic phylogenetic outgroup exists.
Under this regime, \RhizAeon falls back onto outgroup-invariant triplet criteria:
\begin{enumerate}
\item \textbf{Hypergeometric Contingency}: Fisher's exact test across informative sites within versus outside the tract, calibrated against the exact number of unordered candidate triplets:
\begin{equation}
p \le p_{\text{crit}} = \frac{\alpha}{\binom{N}{3}} = \frac{6\alpha}{N(N-1)(N-2)}.
\label{eq:methods_bonferroni_triplets}
\end{equation}
\item \textbf{Poisson Informative Floor}: A minimal count of distinguishing mutations donated by the minor parent ($s_{\text{donor}} \ge 3$).
\item \textbf{Directional Recoil Significance}: Standardized transition displacement along the parental dipole axis exceeding Gaussian extreme-value bounds ($\Delta_{\text{net}} \ge z_{\text{crit}} \cdot \text{SE}$).
\item \textbf{Private Autapomorphy Sieve}: Bounding private autapomorphic mutations along the candidate branch ($\rho_{\text{priv}} < 0.50$) to reject localized heterotachy spikes.
\end{enumerate}
This invariant fallback guarantees rigorous false positive control without requiring external reference taxa or inducing false negatives in closed outbreak cohorts.

\subsection{The Reticulation Fidelity Scoring Framework}
\label{subsec:methods_scoring_framework}

A simulated ground-truth recombination event is formalized as an oriented tuple:
\begin{equation}
E = (C, H, D, \mathcal{I}),
\label{eq:methods_event_tuple}
\end{equation}
where $C \in \{1, \dots, N\}$ denotes the recombinant child isolate, $H \in \{1, \dots, N\}$ denotes the major parental recipient lineage (home), $D \in \{1, \dots, N\} \cup \{\text{Ghost}\}$ denotes the minor parental donor lineage, and $\mathcal{I} \subset [1, L]$ denotes the genomic interval of horizontal genetic exchange.
An inferred event is denoted $\hat{E} = (\hat{C}, \hat{H}, \hat{D}, \hat{\mathcal{I}})$.

Let $\mathcal{S}_{\text{inf}}(H, D) \subset \{1, \dots, L\}$ denote the set of informative polymorphisms distinguishing parental lineages $H$ and $D$:
\begin{equation}
\mathcal{S}_{\text{inf}}(H, D) = \{s \in \{1, \dots, L\} \mid H[s] \neq D[s]\}.
\label{eq:methods_inf_set}
\end{equation}
For true exchange interval $\mathcal{I}$ and inferred interval $\hat{\mathcal{I}}$, the sets of true and inferred transferred informative polymorphisms are:
\begin{equation}
\mathcal{S}_{\text{true}} = \{s \in \mathcal{S}_{\text{inf}}(H, D) \mid s \in \mathcal{I}\}, \qquad \mathcal{S}_{\text{inferred}} = \{s \in \mathcal{S}_{\text{inf}}(H, D) \mid s \in \hat{\mathcal{I}}\}.
\end{equation}
The spatial accuracy score is formulated as the Jaccard similarity over informative sites:
\begin{equation}
J_{\text{SNP}}(\mathcal{I}, \hat{\mathcal{I}}) = \frac{|\mathcal{S}_{\text{true}} \cap \mathcal{S}_{\text{inferred}}|}{|\mathcal{S}_{\text{true}} \cup \mathcal{S}_{\text{inferred}}|}.
\label{eq:methods_j_snp}
\end{equation}
Alignments contain zero phylogenetic signal between adjacent informative sites.
Consequently, the likelihood profile forms an exact mathematical plateau.
Whenever predicted boundaries fall within the flanking uninformative plateaus of the true exchange ($E_{\text{SNP}} = 0$), the informative sets coincide ($\mathcal{S}_{\text{inferred}} \equiv \mathcal{S}_{\text{true}}$), yielding $J_{\text{SNP}} = 1.00$.
The metric naturally scales across evolutionary regimes, penalizing predictions that traverse contradictory mutations without imposing arbitrary nucleotide distance thresholds.

Lineage identification requires genealogical discipline.
To evaluate parental attribution without penalizing close sister isolates or framing innocent taxa, we define genealogical kinship affinity.
Let $d_{\text{tree}}(X, Y)$ denote the patristic distance between taxa $X$ and $Y$ on the reference phylogeny.
The genealogical affinity kernel is defined as:
\begin{equation}
w(X, Y) = \exp\left(-\frac{d_{\text{tree}}(X, Y)}{\sigma_{\text{clade}}}\right),
\label{eq:methods_patristic_kernel}
\end{equation}
where $\sigma_{\text{clade}}$ denotes the characteristic divergence scale of the surrounding monophyletic sublineage.
The directed triad attribution score is formulated as:
\begin{equation}
S_{\text{triad}}(C, H, D;\, \hat{C}, \hat{H}, \hat{D}) = w(C, \hat{C}) \times \left( \frac{w(H, \hat{H}) + w(D, \hat{D})}{2} \right).
\label{eq:methods_s_triad}
\end{equation}
The child affinity $w(C, \hat{C})$ operates as an outer multiplicative gate.
If an algorithm misidentifies an unreticulated background isolate as the recombinant child, $w(C, \hat{C}) \approx 0$, collapsing the triad score.
Symmetrically, if an algorithm inverts the reticulation direction by swapping child and donor ($\hat{C} \approx D$), the deep divergence between recombinant and donor lineages extinguishes $w(C, D)$, purging inverted calls without ad-hoc rules.
When the true donor represents an unsampled ghost lineage ($D = \text{Ghost}$), identifying an unsampled donor awards full credit ($w(D, \hat{D}) = 1.00$ for $\hat{D} = \text{Ghost}$), while naming a sampled sister lineage receives partial credit proportional to its patristic distance to the introgression node.

For each pair of true event $E_i$ and inferred event $\hat{E}_j$, the composite event affinity is the product of lineage fidelity and spatial accuracy:
\begin{equation}
s(E_i, \hat{E}_j) = S_{\text{triad}}(E_i, \hat{E}_j) \times J_{\text{SNP}}(E_i, \hat{E}_j).
\label{eq:methods_event_affinity}
\end{equation}
To evaluate reticulations under strict cardinality discipline, we construct the pairwise cost matrix $C_{ij} = 1 - s(E_i, \hat{E}_j)$ for all $i \in \{1, \dots, K\}$ and $j \in \{1, \dots, \hat{K}\}$.
We compute the optimal bijective bipartite matching $\mathcal{M}$ using the Hungarian algorithm \cite{kuhn1955hungarian}.
Every unmatched true event counts as a false negative, while every unmatched predicted event counts as an ungrounded false alarm.

Across the optimal matching $\mathcal{M}$, reticulation precision and recall are defined as:
\begin{align}
\text{PPV}_{\mathcal{R}} &= \frac{1}{\hat{K}} \sum_{(E_i, \hat{E}_j) \in \mathcal{M}} s(E_i, \hat{E}_j), \label{eq:methods_ppv_r} \\
\text{TPR}_{\mathcal{R}} &= \frac{1}{K} \sum_{(E_i, \hat{E}_j) \in \mathcal{M}} s(E_i, \hat{E}_j). \label{eq:methods_tpr_r}
\end{align}
The harmonic mean of precision and recall defines the composite Reticulation Fidelity Index:
\begin{equation}
\mathcal{R}_{F_1} = \frac{2 \cdot \text{PPV}_{\mathcal{R}} \cdot \text{TPR}_{\mathcal{R}}}{\text{PPV}_{\mathcal{R}} + \text{TPR}_{\mathcal{R}}} = \frac{2 \sum_{(E_i, \hat{E}_j) \in \mathcal{M}} s(E_i, \hat{E}_j)}{K + \hat{K}}.
\label{eq:methods_rf1}
\end{equation}
On non-recombinant clonal null alignments ($K = 0$), the index evaluates to $\mathcal{R}_{F_1} = 1.00$ if $\hat{K} = 0$, and $\mathcal{R}_{F_1} = 0.00$ if any false positive event is declared ($\hat{K} > 0$), with alignment-level false positive rate $\text{FPR} = \mathbb{I}(\hat{K} > 0)$.
To summarize performance across all four biological axes, benchmarks report the multi-dimensional evaluation profile:
\begin{equation}
\mathbf{S}_{\mathcal{R}} = \left(\text{TPR}_{\mathcal{R}},\, \text{PPV}_{\mathcal{R}},\, \overline{S}_{\text{triad}},\, \overline{J}_{\text{SNP}}\right).
\label{eq:methods_score_profile}
\end{equation}

\subsection{Synthetic Benchmark Protocols: Epidemic Clonal Expansions and Multi-Way Mosaics}
\label{subsec:methods_benchmark_protocols}

To rigorously assess algorithmic fidelity under authentic evolutionary challenges without synthetic placeholders or parametric shortcuts, we implemented automated benchmark generators operating on raw sequence alignments under continuous-time Markov substitution dynamics.
All evolutionary simulations were conducted under the HKY85 substitution model with transition/transversion bias $\kappa = 2.5$ and empirical base frequencies $\pi = [0.30, 0.20, 0.20, 0.30]$ using exact matrix exponentiation $P(t) = \exp(Q t)$.

\paragraph{Scenario 1.1: Epidemic Clonal Expansion Suite.}
To assess event consolidation and parental attribution discipline under epidemic lineage growth, we simulated an ancestral recombinant node descending from two parental lineages separated by branch length $d_{\text{parent}} \in \{0.02, 0.05, 0.10\}$ ($L = 5,000$ nt).
The ancestral recombinant acquired a 1,500 nt cassette ($[1,500\text{--}3,000]$) from $P_2$ in a $P_1$ background.
The recombinant node subsequently underwent epidemic transmission expansion, sprouting $K \in \{2, 5, 8, 10, 15\}$ contemporary descendant isolates, each accumulating independent transmission drift ($\delta = 0.002$).
The alignment contained 40 strictly non-recombinant background parental isolates (20 in $P_1$ and 20 in $P_2$), yielding $N = 42\text{--}55$ taxa.
Twenty independent replicates were generated for each of the 15 parameter conditions (300 alignments total).

\paragraph{Scenario 1.2: Multi-Way Mosaic Swarms Suite.}
To evaluate algorithmic limits when the two-parent assumption is violated, we simulated 540 alignments ($L = 6,000$ nt; 20 replicates across 27 parameter conditions) across three distinct multi-parental topologies:
\begin{enumerate}
\item \textbf{Dual-Cassette Swarms}: A $P_1$ background acquiring two independent foreign cassettes from $P_2$ ($[1,200\text{--}2,400]$, length 1,201 nt) and $P_3$ ($[3,600\text{--}4,800]$, length 1,201 nt), leaving $P_1$ on flanks and interstitial spacer.
\item \textbf{Nested Russian Doll Swarms}: A $P_1$ background acquiring an outer introgression from $P_2$ ($[1,500\text{--}4,500]$, length 3,001 nt), within which an internal cassette ($[2,500\text{--}3,500]$, length 1,001 nt) is replaced by $P_3$, modeling secondary crossover within a recombinant lineage.
\item \textbf{Triple Crossover Cascades}: Sequential serial crossovers traversing the three parental lineages ($P_1 [1\text{--}2,000] \to P_2 [2,001\text{--}4,000] \to P_3 [4,001\text{--}6,000]$).
\end{enumerate}
Each topology was evaluated across divergence levels $d_{\text{parent}} \in \{0.02, 0.05, 0.10\}$ and descendant swarm sizes $K \in \{1, 5, 10\}$ with 45 non-recombinant background taxa (15 per parental clade; $N = 46\text{--}55$ taxa).
Performance was evaluated against 3SEQ using identical alignments, measuring background false positive rate ($\text{FPR}_{\text{parent}}$), precision, event cardinality ($\hat{K}_{\text{unique}}$), spatial Jaccard similarity, and execution runtime.
